% Irodov section 2.1, original numbered problems 21--25,
% printed p.65, PDF p.74. Van der Waals equation for one mole:
% p=RT/(V-b)-a/V^2, V is molar volume (one mole).
\thermoproblem{21}{65}
在温度 $T=300\,\mathrm K$ 时，
二氧化碳的质量密度达到
$\rho=500\,\mathrm{g/L}$，
需有多大压强？
分别用理想气体方程和范德瓦耳斯方程计算。
\exsolution
设二氧化碳的摩尔质量
$M=44.0\,\mathrm{g/mol}$，
每摩尔的体积为
\[
V_m=\frac M\rho
=\frac{44.0\,\mathrm{g/mol}}{500\,\mathrm{g/L}}
=0.0880\,\mathrm{L/mol}.
\]
理想气体方程给
\[
p_{\rm id}=\frac{RT}{V_m},
\]
取 $R=0.082057\,\mathrm{L\,atm/(mol\,K)}$，
\[
\boxed{p_{\rm id}\simeq280\,\mathrm{atm}}.
\]
范德瓦耳斯方程为
$(p+a/V_m^2)(V_m-b)=RT$，
解得
\[
\boxed{p_{\rm vdW}=
\frac{RT}{V_m-b}-\frac{a}{V_m^2}}.
\]
采用该气体常用近似参数
$a\simeq3.59\,\mathrm{atm\,L^2/mol^2}$、
$b\simeq0.0427\,\mathrm{L/mol}$，
代入
\[
p_{\rm vdW}
\simeq\frac{0.082057(300)}{0.0880-0.0427}
-\frac{3.59}{(0.0880)^2}
\simeq\boxed{80\,\mathrm{atm}}.
\]
\excheck 两结果差异很大，因为该密度下
分子排除体积 $b$ 已与摩尔体积同量级，
理想气体近似失效。

\thermoproblem{22}{65}
容器体积为 $V=1.00\,\mathrm L$，
内装一摩尔氮气。
求在多高的温度下，
理想气体方程计算的压强与范德瓦耳斯方程
计算值的相对误差为 $\eta=10\%$
（相对范德瓦耳斯压强定义）；
并求此时压强。
\exsolution
取这一摩尔气体的两种压强模型：
\[
p_{\rm id}=\frac{RT}{V},\qquad
p_{\rm vdW}=\frac{RT}{V-b}-\frac a{V^2}.
\]
按本题考虑的分支，理想气体压强高于
范德瓦耳斯值，故
$(p_{\rm id}-p_{\rm vdW})/p_{\rm vdW}=\eta$，
即
$p_{\rm vdW}=p_{\rm id}/(1+\eta)$。
代入并整理：
\[
\frac{RT}{V-b}-\frac a{V^2}
=\frac{RT}{V(1+\eta)}.
\]
移项求温度：
\[
\frac a{V^2}
=RT\frac{\eta V+b}{V(1+\eta)(V-b)},
\]
\[
\boxed{
T=\frac{a(1+\eta)(V-b)}
{R V(\eta V+b)}}.
\]
此时压强
\[
\boxed{p_{\rm vdW}=\frac{RT}{V(1+\eta)}}.
\]
代入氮气的常用近似
$a\simeq1.4\,\mathrm{atm\,L^2/mol^2}$、
$b\simeq0.039\,\mathrm{L/mol}$，
得 $T$ 为约 $1.3\times10^2\,\mathrm K$，
$p_{\rm vdW}$ 为约 $10\,\mathrm{atm}$，
与书末所取参数给出的
$\sim133\,\mathrm K$、
$\sim9.9\,\mathrm{atm}$ 为相同数量级。
\excheck 题面没有提供 $a,b$ 的具体数值，
因此精确数值依赖选用的氮气材料常数表；
不可把简答中的数值视作不依赖 $a,b$ 的普适结果。

\thermoproblem{23}{65}
一摩尔某种气体装在
$V=0.250\,\mathrm L$ 的容器里。
在 $T_1=300\,\mathrm K$ 时
$p_1=90\,\mathrm{atm}$，
在 $T_2=350\,\mathrm K$ 时
$p_2=110\,\mathrm{atm}$。
求该气体的范德瓦耳斯参数 $a,b$。
\exsolution
对同一摩尔体积 $V$，
两状态均满足
\[
p_i=\frac{RT_i}{V-b}-\frac a{V^2}.
\]
相减消去 $a$：
\[
p_2-p_1=\frac{R(T_2-T_1)}{V-b},
\]
所以
\[
\boxed{b=V-\frac{R(T_2-T_1)}{p_2-p_1}}.
\]
再代回任一状态：
\[
\boxed{a=V^2\left(\frac{RT_1}{V-b}-p_1\right)}.
\]
也可代入第一式消去 $b$ 得
\[
a=V^2
\frac{p_1T_2-p_2T_1}{T_1-T_2}.
\]
取 $R=0.082057\,\mathrm{L\,atm/(mol\,K)}$，
数值为
\[
\boxed{b\simeq0.0449\,\mathrm{L/mol}},\qquad
\boxed{a\simeq1.88\,\mathrm{atm\,L^2/mol^2}}.
\]
\excheck 在本版原书简答中，
$a$ 约写成 $1.85$、
$b$ 约为 $0.042\,\mathrm{L/mol}$；
按题面印刷的两组 $p,T$ 和通用 $R$
精确代入得到上式数值。
该差别主要由取整或原版数值表造成，
公式可直接代入反查两个实验压强。

\thermoproblem{24}{65}
一摩尔范德瓦耳斯气体在温度 $T$ 下
摩尔体积为 $V$，其状态方程为
$(p+a/V^2)(V-b)=RT$。
求等温压缩系数
\[
\chi_T=-\frac1V
\left(\frac{\partial V}{\partial p}\right)_T
\]
随 $V$ 的函数。
\exsolution
首先把压强显式写为
\[
p(V,T)=\frac{RT}{V-b}-\frac a{V^2}.
\]
在温度固定时求导：
\[
\left(\frac{\partial p}{\partial V}\right)_T
=-\frac{RT}{(V-b)^2}+\frac{2a}{V^3}.
\]
故
\[
\chi_T=-\frac1V
\frac{1}{(\partial p/\partial V)_T}
=\frac1V
\frac{1}{RT/(V-b)^2-2a/V^3}.
\]
通分得到
\[
\boxed{\chi_T=
\frac{V^2(V-b)^2}
{RTV^3-2a(V-b)^2}}.
\]
\excheck 力学稳定的单相气体须满足
$(\partial p/\partial V)_T<0$，
即上式分母为正；
临界自旋odal附近压缩系数可能显著增大乃至发散。

\thermoproblem{25}{65}
接着上一题，
判断在什么温度条件下，
范德瓦耳斯气体的等温压缩系数
$\chi_{\rm vdW}$ 大于理想气体压缩系数
$\chi_{\rm id}$。
讨论摩尔体积 $V\gg b$ 时的简化结果。
\exsolution
先明确采用\emph{相同温度 $T$ 和相同摩尔体积 $V$}
比较两种状态方程。
理想气体在同一 $V,T$ 时
$p_{\rm id}=RT/V$，
等温压缩系数为
$\chi_{\rm id}=1/p_{\rm id}=V/(RT)$。
范德瓦耳斯气体在力学稳定分支的
压缩系数为
\[
\chi_{\rm vdW}
=\frac1{V[RT/(V-b)^2-2a/V^3]}.
\]
比较
$\chi_{\rm vdW}>\chi_{\rm id}$
且两个分母均为正，等价于
\[
RT>\frac{RTV^2}{(V-b)^2}-\frac{2a}{V}.
\]
移项并使用
$V^2-(V-b)^2=b(2V-b)$，
得到
\[
\frac{2a}{V}>
RT\frac{b(2V-b)}{(V-b)^2},
\]
因此
\[
\boxed{T<
\frac{2a(V-b)^2}{RbV(2V-b)}}.
\]
当 $V\gg b$ 时，
右边趋于 $a/(Rb)$，
从而
\[
\boxed{T<\frac{a}{Rb}\qquad(V\gg b)}.
\]
\excheck \textbf{与本版书末不一致：}
书末只给 $T>a/(Rb)$，方向与按\emph{同 $V,T$}
直接比较压缩系数所得的严格不等式相反。
比较不同气体的“压缩系数”必须明确
是在相同 $T,V$ 还是相同 $T,p$ 条件下；
本稿固定了前者，并将书末差异列入勘误，
不默默改写不等式方向。
